\documentclass[../../main.tex]{subfiles}
\begin{document}

\begin{center}
  {\fontsize{48pt}{58pt}\selectfont \textbf{1963}\normalsize\textbf{年度}}
\end{center}

\begin{table}[htb]
  \centering
  \caption{出題テーマと難易度}
  \begin{tabular}{|c|p{9cm}|c|}
    \hline
    \textbf{問題番号} & \textbf{テーマ} & \textbf{難易度} \\
    \hline
    第1問 & 整数問題 & \\
    \hline
    第2問 & 不等式の証明 & \\
    \hline
    第3問 & 図形と円 & \\
    \hline
    第4問 & 微分（極値） & \\
    \hline
    第5問 & 面積（積分） & \\
    \hline
    第6問 & 積分と証明 & \\
    \hline
  \end{tabular}
\end{table}

\begin{table}[htb]
  \centering
  \caption{解答}
  \renewcommand{\arraystretch}{1.4}
  \begin{tabular}{|c|c|p{9cm}|}
    \hline
    \textbf{問題番号} & \textbf{小問} & \textbf{答え} \\
    \hline
    第1問 & - & $(x,y,z)=(2,1,1)$ \\
    \hline
    第2問 & - & 証明問題 \\
    \hline
    第3問 & (1) & $0<r<2-\sqrt{2}$ \\
    \cline{2-3}
    & (2) & $0<\theta<\dfrac{\pi}{2}$ \\
    \hline
    第4問 & - & 極大 $f(1)=-5$；極小 $f(\sqrt{3})=-6\sqrt{3}+5$ \\
    \hline
    第5問 & - & $S_1:S_2=21:5$ \\
    \hline
    第6問 & - & 証明問題 \\
    \hline
  \end{tabular}
\end{table}

\paragraph{試験形態}
6問．試験時間は要確認．

\end{document}
